% Figure: total annualised cost versus tube-side velocity, decomposed into a
% capital cost that falls with velocity (more h, less area) and a pumping cost
% that rises as velocity cubed, with the optimum at the sum's minimum.
\begin{figure}[htbp]
\centering
\begin{tikzpicture}
\begin{axis}[
  width=12cm, height=8cm,
  xlabel={tube-side velocity $u$ (m\,s$^{-1}$)},
  ylabel={annualised cost (relative)},
  xmin=0.3, xmax=4, ymin=0, ymax=5,
  axis lines=left,
  legend pos=north east,
  legend cell align=left,
  tick label style={font=\scriptsize},
  label style={font=\small},
  legend style={font=\scriptsize},
  samples=80,
  domain=0.4:4,
  clip=false,
]
% capital cost ~ area ~ 1/h ~ u^{-0.8}, scaled
\addplot[engblue, very thick] {2.2*x^(-0.8)};
\addlegendentry{capital (area $\propto 1/h \propto u^{-0.8}$)}
% pumping cost ~ pumping power ~ u^3, scaled small
\addplot[engamber, very thick] {0.18*x^3};
\addlegendentry{pumping ($\propto u^3$)}
% total
\addplot[engred, very thick] {2.2*x^(-0.8) + 0.18*x^3};
\addlegendentry{total}
% optimum marker near u ~ 1.6 (where d/du of total = 0)
\addplot[only marks, mark=*, engred] coordinates {(1.6,2.27)};
\node[engred, font=\scriptsize, anchor=south west] at (axis cs:1.65,2.35) {optimum};
\end{axis}
\end{tikzpicture}
\caption[Heat-transfer versus pressure-drop cost trade-off]{The central
optimisation of exchanger design, expressed as annualised cost against tube-side
velocity. Raising the velocity raises the heat-transfer coefficient roughly as
$u^{0.8}$, so the area, and with it the capital cost, falls. The same velocity
raises the pressure drop as $u^2$ and the pumping power as $u^3$, so the
operating cost climbs steeply. The sum has a flat minimum, and the optimum
velocity for liquids in tubes lands in the range of one to three metres per
second for most services. The minimum is broad, which is fortunate: a design that
sits anywhere near the bottom of the total curve is close to optimal, so the
engineer rarely needs to locate the exact minimum.}
\label{fig:vol04ch06:pressure-drop-tradeoff}
\end{figure}
